% Standalone research note for the hybrid-local-solver project.
%
% This note audits the COLT 2023 ASPR theorem and proves a path-family lower
% bound for the literal nested active-set algorithm. Source claims, repaired
% statements, new proofs, and open problem-level lower bounds are separated.

\documentclass[11pt]{article}

% Common class-neutral preamble for standalone notes under manuscript/notes/.
% Note entry points all live exactly two directories below manuscript/.

\usepackage[margin=1in]{geometry}
\usepackage{amsmath,amssymb,amsthm,mathtools,bm}
\usepackage{algorithm}
\usepackage{algorithmic}
\usepackage{booktabs,tabularx,multirow,array,longtable}
\usepackage{enumitem}
\usepackage{microtype}
\usepackage[numbers,sort&compress]{natbib}
\usepackage{xcolor}
\usepackage[colorlinks=true,citecolor=blue,linkcolor=black,urlcolor=blue]{hyperref}
\usepackage[capitalize,nameinlink,noabbrev]{cleveref}

\newtheorem{theorem}{Theorem}[section]
\newtheorem{lemma}[theorem]{Lemma}
\newtheorem{proposition}[theorem]{Proposition}
\newtheorem{corollary}[theorem]{Corollary}
\newtheorem{conjecture}[theorem]{Conjecture}
\theoremstyle{definition}
\newtheorem{definition}[theorem]{Definition}
\newtheorem{assumption}[theorem]{Assumption}
\newtheorem{problem}[theorem]{Problem}
\newtheorem{observation}[theorem]{Observation}
\theoremstyle{remark}
\newtheorem{remark}[theorem]{Remark}
\theoremstyle{plain}

% Canonical mathematical notation for the active manuscript.
%
% This file preserves the recurring notation style used in the author's
% NeurIPS 2024 and 2025 papers while excluding unrelated tensor and
% machine-learning boilerplate from those archived command collections.
%
% Prefix convention:
%   \v...  bold vectors
%   \m...  bold matrices
%   \g...  calligraphic sets, graphs, and families
%   \s...  blackboard-bold spaces and number systems

\newcommand{\method}{\textsc{HybridLocalSolver}}

% Font and scalar shorthands retained from the archived manuscripts.
\newcommand{\mc}{\mathcal}
\newcommand{\R}{\mathbb{R}}
\newcommand{\E}{\mathbb{E}}
\newcommand{\G}{\mathbb{G}}
\newcommand{\N}{\mathcal{N}}
% Accuracy namespaces are semantic: do not add a bare \eps alias.
% Degree-normalized residual tolerance of classical APPR is kept distinct from
% the objective-gap and proximal fixed-point tolerances of the RPPR sections.
\newcommand{\epsappr}{\varepsilon_{\mathrm{appr}}}
\newcommand{\epsobj}{\varepsilon_{\mathrm{obj}}}
\newcommand{\epspg}{\varepsilon_{\mathrm{pg}}}
\newcommand{\epsppr}{\varepsilon_{\mathrm{ppr}}}
\newcommand{\epsin}{\varepsilon_{\mathrm{in}}}
\newcommand{\epskkt}{\varepsilon_{\mathrm{kkt}}}
\newcommand{\epsburn}{\varepsilon_{\mathrm{burn}}}
\newcommand{\epssol}{\varepsilon_{\mathrm{sol}}}
\newcommand{\epspprinst}[1]{\varepsilon_{\mathrm{ppr},#1}}
\newcommand{\trans}{\mathsf{T}}
\newcommand{\grad}{\nabla}

% Delimiters and generic helpers.
% Norms and inner products use fixed-size delimiters by default.
% Use an optional size (e.g. \norm[\big]{...}) or * for automatic sizing.
\DeclarePairedDelimiter{\norm}{\lVert}{\rVert}
\newcommand{\abs}[1]{\left\lvert #1 \right\rvert}
\DeclarePairedDelimiterX{\inner}[2]{\langle}{\rangle}{#1, #2}
\newcommand{\ip}[2]{\inner{#1}{#2}}
\newcommand{\card}[1]{\left\lvert #1 \right\rvert}
\newcommand{\set}[1]{\left\{#1\right\}}
\newcommand{\ceil}[1]{\left\lceil #1 \right\rceil}
\newcommand{\floor}[1]{\left\lfloor #1 \right\rfloor}
\newcommand{\comp}[1]{\overline{#1}}
\newcommand{\conj}[1]{\overline{#1}}
\newcommand{\mean}[1]{\overline{#1}}
\newcommand{\bigO}{\mathcal{O}}
\newcommand{\softO}{\widetilde{\mathcal{O}}}
\newcommand{\softOmega}{\widetilde{\Omega}}
\newcommand{\pos}[1]{\left[#1\right]_{+}}

% Bold lowercase vectors.
\newcommand{\va}{\bm{a}}
\newcommand{\vb}{\bm{b}}
\newcommand{\vc}{\bm{c}}
\newcommand{\vd}{\bm{d}}
\newcommand{\ve}{\bm{e}}
\newcommand{\vf}{\bm{f}}
\newcommand{\vg}{\bm{g}}
\newcommand{\vh}{\bm{h}}
\newcommand{\vi}{\bm{i}}
\newcommand{\vj}{\bm{j}}
\newcommand{\vk}{\bm{k}}
\newcommand{\vl}{\bm{l}}
\newcommand{\vm}{\bm{m}}
\newcommand{\vn}{\bm{n}}
\newcommand{\vo}{\bm{o}}
\newcommand{\vp}{\bm{p}}
\newcommand{\vq}{\bm{q}}
\newcommand{\vr}{\bm{r}}
\newcommand{\vs}{\bm{s}}
\newcommand{\vt}{\bm{t}}
\newcommand{\vu}{\bm{u}}
\newcommand{\vv}{\bm{v}}
\newcommand{\vw}{\bm{w}}
\newcommand{\vx}{\bm{x}}
\newcommand{\vy}{\bm{y}}
\newcommand{\vz}{\bm{z}}

% Bold Greek vectors and special vectors.
\newcommand{\valpha}{\bm{\alpha}}
\newcommand{\vbeta}{\bm{\beta}}
\newcommand{\vdelta}{\bm{\delta}}
\newcommand{\vepsilon}{\bm{\epsilon}}
\newcommand{\veta}{\bm{\eta}}
\newcommand{\vgamma}{\bm{\gamma}}
\newcommand{\vlambda}{\bm{\lambda}}
\newcommand{\vmu}{\bm{\mu}}
\newcommand{\vomega}{\bm{\omega}}
\newcommand{\vphi}{\bm{\phi}}
\newcommand{\vpi}{\bm{\pi}}
\newcommand{\vpsi}{\bm{\psi}}
\newcommand{\vrho}{\bm{\rho}}
\newcommand{\vsigma}{\bm{\sigma}}
\newcommand{\vtheta}{\bm{\theta}}
\newcommand{\vxi}{\bm{\xi}}
\newcommand{\vell}{\bm{\ell}}
\newcommand{\vDelta}{\bm{\Delta}}
\newcommand{\vGamma}{\bm{\Gamma}}
\newcommand{\vLambda}{\bm{\Lambda}}
\newcommand{\vPhi}{\bm{\Phi}}
\newcommand{\vSigma}{\bm{\Sigma}}
\newcommand{\vzero}{\bm{0}}
\newcommand{\vone}{\bm{1}}
\newcommand{\one}{\mathbf{1}}
\newcommand{\zero}{\vzero}
\newcommand{\eunit}[1]{\bm{e}_{#1}}
\newcommand{\vDeltax}{\bm{\Delta}\bm{x}}

% Bold uppercase matrices.
\newcommand{\mA}{\bm{A}}
\newcommand{\mB}{\bm{B}}
\newcommand{\mC}{\bm{C}}
\newcommand{\mD}{\bm{D}}
\newcommand{\mE}{\bm{E}}
\newcommand{\mF}{\bm{F}}
\newcommand{\mG}{\bm{G}}
\newcommand{\mH}{\bm{H}}
\newcommand{\mI}{\bm{I}}
\newcommand{\mJ}{\bm{J}}
\newcommand{\mK}{\bm{K}}
\newcommand{\mL}{\bm{L}}
\newcommand{\mM}{\bm{M}}
\newcommand{\mN}{\bm{N}}
\newcommand{\mO}{\bm{O}}
\newcommand{\mP}{\bm{P}}
\newcommand{\mQ}{\bm{Q}}
\newcommand{\mR}{\bm{R}}
\newcommand{\mS}{\bm{S}}
\newcommand{\mT}{\bm{T}}
\newcommand{\mU}{\bm{U}}
\newcommand{\mV}{\bm{V}}
\newcommand{\mW}{\bm{W}}
\newcommand{\mX}{\bm{X}}
\newcommand{\mY}{\bm{Y}}
\newcommand{\mZ}{\bm{Z}}

% Bold Greek matrices retained from the graph papers.
\newcommand{\mBeta}{\bm{\beta}}
\newcommand{\mDelta}{\bm{\Delta}}
\newcommand{\mGamma}{\bm{\Gamma}}
\newcommand{\mLambda}{\bm{\Lambda}}
\newcommand{\mOmega}{\bm{\Omega}}
\newcommand{\mPhi}{\bm{\Phi}}
\newcommand{\mPi}{\bm{\Pi}}
\newcommand{\mSigma}{\bm{\Sigma}}
\newcommand{\mTheta}{\bm{\Theta}}

% Calligraphic symbols.
\newcommand{\gA}{\mathcal{A}}
\newcommand{\gB}{\mathcal{B}}
\newcommand{\gC}{\mathcal{C}}
\newcommand{\gD}{\mathcal{D}}
\newcommand{\gE}{\mathcal{E}}
\newcommand{\gF}{\mathcal{F}}
\newcommand{\gG}{\mathcal{G}}
\newcommand{\gH}{\mathcal{H}}
\newcommand{\gI}{\mathcal{I}}
\newcommand{\gJ}{\mathcal{J}}
\newcommand{\gK}{\mathcal{K}}
\newcommand{\gL}{\mathcal{L}}
\newcommand{\gM}{\mathcal{M}}
\newcommand{\gN}{\mathcal{N}}
\newcommand{\gO}{\mathcal{O}}
\newcommand{\gP}{\mathcal{P}}
\newcommand{\gQ}{\mathcal{Q}}
\newcommand{\gR}{\mathcal{R}}
\newcommand{\gS}{\mathcal{S}}
\newcommand{\gT}{\mathcal{T}}
\newcommand{\gU}{\mathcal{U}}
\newcommand{\gV}{\mathcal{V}}
\newcommand{\gW}{\mathcal{W}}
\newcommand{\gX}{\mathcal{X}}
\newcommand{\gY}{\mathcal{Y}}
\newcommand{\gZ}{\mathcal{Z}}

% Blackboard-bold symbols.
\newcommand{\sA}{\mathbb{A}}
\newcommand{\sB}{\mathbb{B}}
\newcommand{\sC}{\mathbb{C}}
\newcommand{\sD}{\mathbb{D}}
\newcommand{\sE}{\mathbb{E}}
\newcommand{\sF}{\mathbb{F}}
\newcommand{\sG}{\mathbb{G}}
\newcommand{\sH}{\mathbb{H}}
\newcommand{\sI}{\mathbb{I}}
\newcommand{\sJ}{\mathbb{J}}
\newcommand{\sK}{\mathbb{K}}
\newcommand{\sL}{\mathbb{L}}
\newcommand{\sM}{\mathbb{M}}
\newcommand{\sN}{\mathbb{N}}
\newcommand{\sO}{\mathbb{O}}
\newcommand{\sP}{\mathbb{P}}
\newcommand{\sQ}{\mathbb{Q}}
\newcommand{\sR}{\mathbb{R}}
\newcommand{\sS}{\mathbb{S}}
\newcommand{\sT}{\mathbb{T}}
\newcommand{\sU}{\mathbb{U}}
\newcommand{\sV}{\mathbb{V}}
\newcommand{\sW}{\mathbb{W}}
\newcommand{\sX}{\mathbb{X}}
\newcommand{\sY}{\mathbb{Y}}
\newcommand{\sZ}{\mathbb{Z}}

% Frequently used operators.
\DeclareMathOperator{\Cov}{Cov}
\DeclareMathOperator{\Var}{Var}
\DeclareMathOperator{\diag}{diag}
\DeclareMathOperator{\nei}{Nei}
\DeclareMathOperator{\nnz}{nnz}
\DeclareMathOperator{\pr}{pr}
\DeclareMathOperator{\prox}{prox}
\DeclareMathOperator{\push}{push}
\DeclareMathOperator{\res}{res}
\DeclareMathOperator{\rel}{Rel}
\DeclareMathOperator{\relu}{ReLU}
\DeclareMathOperator{\sign}{sign}
\DeclareMathOperator{\supp}{supp}
\DeclareMathOperator{\tr}{tr}
\DeclareMathOperator{\Tr}{Tr}
\DeclareMathOperator{\vol}{vol}
\DeclareMathOperator{\work}{work}
\DeclareMathOperator{\Work}{Work}
\DeclareMathOperator*{\argmax}{arg\,max}
\DeclareMathOperator*{\argmin}{arg\,min}

% Project-wide names and claim-status labels used by standalone research notes.
% Mathematical notation belongs in math_commands.tex.

% Algorithms and solver families.
\newcommand{\AESP}{\textsc{AESP}}
\newcommand{\APPR}{\textsc{APPR}}
\newcommand{\ASPR}{\textsc{ASPR}}
\newcommand{\FISTA}{\textsc{FISTA}}
\newcommand{\ISTA}{\textsc{ISTA}}
\newcommand{\LOCAPPR}{\textsc{LocAPPR}}
\newcommand{\LOCGD}{\textsc{LocGD}}
\newcommand{\LOCSOR}{\textsc{LocSOR}}
\newcommand{\LOCCH}{\textsc{LocCH}}
\newcommand{\LOCHB}{\textsc{LocHB}}
\newcommand{\LocProxGS}{\textsc{LocProxGS}}
\newcommand{\Hybrid}{\textsc{AESP--LocSOR}}

% Claim classifications required by the repository research-note protocol.
\newcommand{\statussource}{\textbf{Source result}}
\newcommand{\statusproved}{\textbf{Proved here}}
\newcommand{\statusconditional}{\textbf{Conditional}}
\newcommand{\statusconjecture}{\textbf{Open / conjectural}}
\newcommand{\statusopen}{\textbf{Open}}
\newcommand{\statusempirical}{\textbf{Empirical}}
\newcommand{\statusobservation}{\textbf{Empirical observation}}
\newcommand{\statusrefuted}{\textbf{Refuted route}}

% Compatibility names used by the older complete-note presentation layer.
\newcommand{\Status}[2]{\textbf{#1:} #2}
\newcommand{\SourceStatus}{\textnormal{\textsc{Source result}}}
\newcommand{\ProvedStatus}{\textnormal{\textsc{Proved here}}}
\newcommand{\ConditionalStatus}{\textnormal{\textsc{Conditional}}}
\newcommand{\ComputedStatus}{\textnormal{\textsc{Computed}}}
\newcommand{\RefutedStatus}{\textnormal{\textsc{Refuted route}}}
\newcommand{\OpenStatus}{\textnormal{\textsc{Open}}}



\title{The ASPR 2023 Bound Revisited:\\
Correctness, Repeated Restricted Solves, and a Tight Path Lower Bound}
\author{Baojian Zhou\\
\small Working research note for \texttt{hybrid-local-solver}}
\date{August 2026}

\begin{document}
\maketitle

\begin{abstract}
Mart\'inez-Rubio, Wirth, and Pokutta proposed the accelerated sparse PageRank
method (\ASPR{}), which repeatedly solves an expanding restricted
positive-orthant quadratic by accelerated projected gradient descent, retracts
the approximate restricted minimizer, and adds every negative-gradient
coordinate.  Its stated work is
\[
 \softO\!\left(
   |S^\star|\,\widetilde{\vol}(S^\star)\sqrt{L/\alpha}
   +|S^\star|\vol(S^\star)
 \right).
\]
We first audit the proof.  The intended support-safety and objective-gap
argument is correct after repairing the quadratic normalization, the printed
RPPR linear term, and one distance display that uses the retracted point where
the APGD output is required.  Several further defects affect constants,
notation, or the executable access model but not the repaired theorem.
We also audit the authors' later Julia implementation.  Its official plots do
not show default \ASPR{} to be slow relative to its FISTA baseline; they
usually show \ASPR{} faster and CASPR fastest.  However, the optional
periodic-gradient variants delete the active-to-boundary matrix block before
their purported full-gradient check, so early discovery cannot fire on RPPR
and the variants only add work.  The in-place retraction also fails to clip
newly negative entries.

We then prove that the factor $|S^\star|$ is genuine.  On an endpoint-seeded
RPPR path whose regularized solution is positive on all $s$ vertices, every
proper active prefix exposes exactly its next vertex and no other vertex.
For all sufficiently small objective tolerances, the retraction preserves this
strict frontier gradient, so literal \ASPR{} makes exactly $s$ restricted
APGD calls.  The prescribed inner count is
$\Omega(1/\sqrt\alpha)$ at every stage, while a restricted gradient on a
prefix of length $r$ costs $\Omega(r)$.  Consequently
\[
 \Work_{\ASPR}
 =\Omega\!\left(\frac{s^2}{\sqrt\alpha}\right),
\]
and literal full-gradient discovery additionally costs $\Omega(s^2)$.
Because both internal and incident path volume are $\Theta(s)$, this matches
the leading published upper bound up to logarithms.  Standard FISTA on the
same finite path avoids the repeated prefix solves.  Quantitatively, taking
$\alpha=s^{-2}$, $\rho=\alpha/100$, and
$\varepsilon_{\mathrm{obj}}=10^{-4}\alpha^2$ gives
$\Omega(s^3\log s)$ work for literal \ASPR{} but only
$O(s^2\log s)$ for full-space FISTA.  Thus restarting acceleration after
every safe expansion causes an actual factor-$s$ loss on this family.  These
are algorithm-specific results, not an oracle lower bound for every local
PageRank method.
\end{abstract}

\tableofcontents

\section{Scope, status, and source map}
\label{sec:scope}

This note concerns Algorithm~4 and Theorems~7--8 of
\citet{martinezrubio2023}.  It does not audit the separate conjugate-directions
method except where the two algorithms share geometric preliminaries.

\begin{center}
\begin{tabularx}{\textwidth}{@{}p{0.19\textwidth}X@{}}
\toprule
Status & Statement \\
\midrule
\statussource & The nested active-set algorithm, APGD subroutine, retraction,
support-safety theorem, and stated work bound are from
\citet{martinezrubio2023}. \\
\statusproved & After the repairs in \cref{sec:audit}, the exact-arithmetic
support-containment and objective-gap argument is valid for strongly convex
quadratics with Stieltjes Hessian. \\
\statusproved & An endpoint-seeded RPPR path with full optimal support forces
one new coordinate per outer stage.  For sufficiently small objective
tolerance, literal \ASPR{} uses
$\Omega(s^2/\sqrt\alpha)$ restricted-solve work and
$\Omega(s^2)$ discovery work. \\
\statusproved & With $\alpha=s^{-2}$, $\rho=\alpha/100$, and
$\varepsilon_{\rm obj}=10^{-4}\alpha^2$, the same-tolerance work is
$\Omega(s^3\log s)$ for literal \ASPR{} and $O(s^2\log s)$ for
full-space \FISTA{}, an actual factor-$s$ separation. \\
\statusconditional & The published $O(|S^\star|)$ space claim is executable
with sorted adjacency lists and a streaming multiway merge, at an additional
logarithmic time factor.  A hash-accumulator implementation instead uses
boundary-sized memory. \\
\statusproved & In official commit \texttt{3a169eb}, the periodic
full-gradient option removes the active--boundary cross block before
evaluation, making early discovery inert for RPPR.  The in-place retraction
also fails when an entry lies strictly between zero and the retraction size. \\
\statusempirical & The later official repository plots generally show default
\ASPR{} faster than the included FISTA and ISTA baselines, with CASPR fastest.
They use one timed run and a minimum-degree seed; they are not experiments in
the COLT/arXiv-v1 PDF. \\
\statusopen & A lower bound for every local first-order method, or a matching
lower bound under a fixed repository-wide PageRank residual, is not proved.
\qquad The repository-wide residual convention remains open. \\
\bottomrule
\end{tabularx}
\end{center}

\paragraph{Source pointers.}
The model and RPPR specialization are on PDF pages 4--5 of
\citet{martinezrubio2023}; Proposition~2 is on pages 6--7; Algorithm~3 and
Proposition~5 are on pages 9--10; Algorithm~4 and Theorems~7--8 are on pages
10--11; the correctness proof is on pages 21--23; and the complexity proof is
on page 24.  The COLT/arXiv-v1 source contains no experiments.  A later
official Julia repository, labeled as containing experiments for a journal
version, is available at
\url{https://github.com/ZIB-IOL/AAPPR.jl}.  The classical FISTA work comparison
and its locality limitations are taken from
\citet{fountoulakis2026}, especially PDF pages 3--6 and 18--23.

\paragraph{Accuracy and work boundary.}
Throughout this note, $\varepsilon_{\rm obj}$ is an objective-gap target.  It
is not APPR's degree-normalized threshold, a PPR infinity-error target, or the
proximal fixed-point residual used in the experiments of
\citet{fountoulakis2026}.  A restricted first-order step is charged the
number of stored nonzeros touched in the restricted Hessian; on a path prefix
this is $\Theta(r)$.  A literal expansion scan is charged the incident
nonzeros of the current support.

\section{Source-aligned RPPR and the corrected obstacle problem}
\label{sec:model}

% Canonical source-aligned PageRank/RPPR reference formulation.
%
% This fragment is intentionally heading-free at the section level so that the
% active paper and every standalone note can input it. It is a manuscript
% reference convention, not an implementation-wide residual decision.
% Primary source: Fountoulakis and Martinez-Rubio, arXiv:2602.21138v2,
% PDF p. 3, Section 3 and equation (RPPR).

\subsection{Common source-aligned PageRank and RPPR model}
\label{subsec:shared-source-aligned-problem}

All documents in this manuscript workspace use the following reference
language. It follows the plain italic notation of the comparison paper:
\(x,s,A,D,Q,I\) denote column vectors or matrices as determined by their
displayed domains. This is the scoped exception to the project's usual
bold-vector and bold-matrix typography. A note may impose a stronger
assumption after this block, but it must not redefine these objects.

Let \(G=(V,E)\) be a finite, simple, undirected, unweighted graph without
isolated vertices, let \(V=[n]:=\{1,\ldots,n\}\), and let
\(A\in\{0,1\}^{n\times n}\) be its symmetric adjacency matrix. For
\(i\in V\), write
\[
    \mathcal N(i):=\{j\in V:\{i,j\}\in E\},
    \qquad
    d_i:=|\mathcal N(i)|,
    \qquad
    D:=\diag(d_1,\ldots,d_n).
\]
For \(S\subseteq V\), define
\[
    \mathcal N(S):=\bigcup_{i\in S}\mathcal N(i),
    \qquad
    \partial S:=\mathcal N(S)\setminus S,
    \qquad
    \operatorname{Ext}(S):=V\setminus(S\cup\partial S),
\]
\begin{equation}
    \vol(S):=\sum_{i\in S}d_i.
    \label{eq:shared-volume}
\end{equation}
The support of a vector is
\(\supp(x):=\{i\in[n]:x_i\neq0\}\). Matrix inequalities use the Loewner
order, and all unqualified vector inequalities are coordinatewise.

The seed is a column distribution
\begin{equation}
    s\in\R_+^n,
    \qquad
    \inner{\one}{s}=1.
    \label{eq:shared-seed}
\end{equation}
The single-seed specialization is always written \(s=e_v\), where \(v\in V\)
is the seed vertex. Thus \(s\) is never used as a vertex label. For
\(\alpha\in(0,1]\), define
\begin{equation}
\begin{aligned}
    \mathcal L&:=I-D^{-1/2}AD^{-1/2},\\
    Q&:=\alpha I+\frac{1-\alpha}{2}\mathcal L
      =\frac{1+\alpha}{2}I
       -\frac{1-\alpha}{2}D^{-1/2}AD^{-1/2},\\
    b&:=\alpha D^{-1/2}s.
\end{aligned}
    \label{eq:shared-pagerank-matrices}
\end{equation}
Since \(0\preceq\mathcal L\preceq2I\),
\(\alpha I\preceq Q\preceq I\). The shared smooth PageRank quadratic is
\begin{equation}
    f(x):=\frac12\inner{x}{Qx}-\inner{b}{x},
    \qquad
    \nabla f(x)=Qx-b.
    \label{eq:shared-pagerank-objective}
\end{equation}
Its unique unregularized minimizer and associated lazy personalized PageRank
vector are
\begin{equation}
    x^0:=Q^{-1}b,
    \qquad
    \pi:=D^{1/2}x^0.
    \label{eq:shared-ppr-solution}
\end{equation}

For a regularization parameter \(\rho>0\), reserve \(g_\rho\) for the
nonsmooth weighted \(\ell_1\) term and define
\begin{equation}
    g_\rho(x):=\alpha\rho\norm{D^{1/2}x}_1,
    \qquad
    F_\rho(x):=f(x)+g_\rho(x),
    \qquad
    x^\star(\rho):=\argmin_{x\in\R^n}F_\rho(x).
    \label{eq:shared-rppr-objective}
\end{equation}
When \(\rho\) is fixed, \(x^\star\) abbreviates \(x^\star(\rho)\), never the
unregularized point \(x^0\). Write
\[
    S^\star(\rho):=\supp(x^\star(\rho)),
    \qquad
    I^\star(\rho):=[n]\setminus S^\star(\rho).
\]
The source records \(x^\star(\rho)\geq0\),
\(\vol(S^\star(\rho))\leq1/\rho\), and the coordinatewise conditions
\begin{equation}
    \nabla_i f(x^\star(\rho))
    \in
    \begin{cases}
        \{-\alpha\rho\sqrt{d_i}\}, & x_i^\star(\rho)>0,\\
        [-\alpha\rho\sqrt{d_i},0], & x_i^\star(\rho)=0.
    \end{cases}
    \label{eq:shared-rppr-kkt}
\end{equation}

The common computational unit is one repeated adjacency-list scan: scanning
coordinate \(i\) costs \(d_i\), and scanning a set \(S\) costs \(\vol(S)\).
Iteration count and wall-clock time do not replace this degree-weighted work
measure.

Finally, accuracy symbols remain distinct. The APPR threshold is
\(\epsappr\); \(\epsobj\) denotes an objective-gap target;
\(\epspg\) denotes the source experiment's proximal fixed-point tolerance; and
\(\epsppr\) may denote a document-scoped degree-normalized PPR error target.
No equality or conversion among these quantities is assumed. In particular,
this shared model does not resolve the repository-wide residual decision.


The shared RPPR minimizer is nonnegative.  On the nonnegative orthant the
weighted $\ell_1$ term is linear, so the source-compatible obstacle form is
\begin{equation}
 \min_{x\geq0}G_\rho(x),
 \qquad
 G_\rho(x)
 :=\frac12x^\top Qx-c_\rho^\top x,
 \qquad
 c_\rho:=\alpha D^{-1/2}s-\alpha\rho D^{1/2}\one.
 \label{eq:obstacle}
\end{equation}
Thus
\begin{equation}
 \nabla G_\rho(x)=Qx-c_\rho,
 \qquad
 \alpha I\preceq Q\preceq I,
 \qquad
 Q_{ij}\leq0\quad(i\neq j).
 \label{eq:obstacle-properties}
\end{equation}
The last property makes $Q$ a symmetric nonsingular Stieltjes matrix.  It is
the order structure used by \ASPR{}.

For the proof audit it is useful to state the generic version.  Let
\begin{equation}
 H=H^\top,
 \qquad
 \mu I\preceq H\preceq LI,
 \qquad
 H_{ij}\leq0\quad(i\neq j),
 \label{eq:generic-H}
\end{equation}
and consider
\begin{equation}
 \min_{x\geq0}G(x),
 \qquad G(x):=\frac12x^\top Hx-c^\top x.
 \label{eq:generic-QP}
\end{equation}
The RPPR specialization is $H=Q$, $\mu=\alpha$, and $L=1$.

\section{A corrected exact-arithmetic ASPR contract}
\label{sec:contract}

Let $S$ be the current active set and let
\begin{equation}
 z^S:=\argmin\{G(z):z\geq0,\ \supp(z)\subseteq S\}
 \label{eq:restricted-minimizer}
\end{equation}
be the exact restricted minimizer.  At an outer stage with $r=|S|$, the
source algorithm chooses
\begin{equation}
 \delta_S^2:=\frac{\varepsilon_{\rm obj}\mu}{(1+r)L^2},
 \qquad
 \widehat\varepsilon_S:=\frac{\mu\delta_S^2}{2}
 =\frac{\varepsilon_{\rm obj}\mu^2}{2(1+r)L^2}.
 \label{eq:inner-tolerances}
\end{equation}
It runs projected APGD until its output $\bar x^S$ satisfies
\begin{equation}
 G(\bar x^S)-G(z^S)\leq\widehat\varepsilon_S,
 \label{eq:inner-gap}
\end{equation}
then retracts
\begin{equation}
 x^+:=[\bar x^S-\delta_S\one]_+
 \label{eq:retraction}
\end{equation}
on $S$, leaves all other coordinates zero, and sets
\begin{equation}
 S^+:=S\cup\{i:\nabla_iG(x^+)<0\}.
 \label{eq:expansion}
\end{equation}
It terminates when $S^+=S$.

The executable inner count should be written
\begin{equation}
 K_S
 :=1+\max\left\{0,
 \left\lceil
  2\sqrt{\frac L\mu}\,
  \log\!\left(
   \frac{(L-\mu)\|\nabla_SG(x)\|_2^2}
        {2\widehat\varepsilon_S\mu^2}
  \right)
 \right\rceil\right\}.
 \label{eq:corrected-inner-count}
\end{equation}
The positive-part clamp is absent from the source pseudocode but is necessary
when the starting point is already sufficiently accurate or has zero
restricted gradient.  It changes no upper bound.

\begin{lemma}[Retraction envelope and restricted gap]
\label{lem:retraction}
Under \cref{eq:inner-gap}, assume also that $z_i^S>0$ for every $i\in S$.
This is the invariant maintained by the ASPR expansion rule.
\begin{equation}
 \|\bar x^S-z^S\|_2\leq\delta_S,
 \qquad
 0\leq x^+\leq z^S,
 \qquad
 G(x^+)-G(z^S)\leq\frac{\varepsilon_{\rm obj}\mu}{L}.
 \label{eq:retraction-properties}
\end{equation}
\end{lemma}

\begin{proof}
Strong convexity and \cref{eq:inner-gap} give
\[
 \|\bar x^S-z^S\|_2^2
 \leq\frac{2\widehat\varepsilon_S}{\mu}
 =\delta_S^2.
\]
If $x_i^+>z_i^S$, then
$\bar x_i^S-z_i^S>\delta_S$, contradicting this distance bound.  Hence
$x^+\leq z^S$.  Since both points lie in the same face and the restricted
gradient of $G$ vanishes at the positive restricted minimizer, smoothness and
nonexpansiveness of projection give
\begin{align*}
 G(x^+)-G(z^S)
 &\leq\frac L2\|x^+-z^S\|_2^2\\
 &\leq L\bigl(\|\bar x^S-z^S\|_2^2+r\delta_S^2\bigr)\\
 &\leq L(1+r)\delta_S^2
 =\frac{\varepsilon_{\rm obj}\mu}{L}.
\end{align*}
\end{proof}

\begin{lemma}[Safe discovery and progress alternative]
\label{lem:safe-progress}
Assume the current active set is contained in the global optimal support.
Every newly discovered coordinate in \cref{eq:expansion} also belongs to that
support.  Moreover, either
\begin{equation}
 G(x^+)-G(x^\star)\leq\varepsilon_{\rm obj}
 \label{eq:accurate-alternative}
\end{equation}
or $S^+\supsetneq S$.
\end{lemma}

\begin{proof}
The envelope $x^+\leq z^S$ and the nonpositive off-diagonal entries imply
that, for $j\notin S$,
\[
 \nabla_jG(x^+)\geq\nabla_jG(z^S).
\]
Thus a negative gradient at $x^+$ is also negative at $z^S$.  Enlarging the
restricted problem by that coordinate makes it strictly positive, and
inverse-positivity of principal Stieltjes matrices shows that the enlarged
restricted minimizer is dominated by $x^\star$.  The coordinate therefore
belongs to the global support.

For the second claim, let
\[
 T(x^+):=\Pi_{\mathbb R_+^n}
 \left(x^+-\frac1L\nabla G(x^+)\right).
\]
Projected gradient descent for a $\mu$-strongly convex, $L$-smooth objective
satisfies
\[
 G(T(x^+))-G(x^\star)
 \leq\left(1-\frac\mu L\right)
       (G(x^+)-G(x^\star)).
\]
If the global gap exceeds $\varepsilon_{\rm obj}$, the progress of this step
is strictly larger than $\varepsilon_{\rm obj}\mu/L$, which is at least the
restricted gap in \cref{eq:retraction-properties}.  Hence
$G(T(x^+))<G(z^S)$.  The point $T(x^+)$ cannot lie in the current face, where
$z^S$ is optimal.  Since $x_j^+=0$ outside $S$, an outside coordinate of
$T(x^+)$ is positive exactly when its gradient at $x^+$ is negative.  Thus
the active set expands.
\end{proof}

Lemmas~\ref{lem:retraction}--\ref{lem:safe-progress} give the intended
Theorem~7: all active sets stay inside $S^\star$, and a stable stage is
$\varepsilon_{\rm obj}$-optimal.

\section{Audit of the published proof}
\label{sec:audit}

The repaired argument above is short, but it is not literally the argument
printed in every display.  The defects fall into three classes.

\subsection{Essential repairs}

\begin{enumerate}[leftmargin=*]
\item The generic source model writes $\langle x,Qx\rangle-\langle b,x\rangle$
while calling $Q$ its Hessian.  The Hessian is $2Q$.  Every subsequent
gradient, smoothness, and PageRank formula instead uses
$\frac12x^\top Qx-b^\top x$.  The latter is the only consistent
interpretation.

\item Source Equation~(4) gives an incorrect expansion of the RPPR obstacle
objective.  The correct linear term is
\[
 \alpha\bigl(\rho D^{1/2}\one-D^{-1/2}s\bigr)^\top x,
\]
as in \cref{eq:obstacle}.  The sign of the seed term and the placement of
$D^{1/2}$ on the penalty are both wrong in the printed expansion.

\item In the proof of Theorem~7, the display that is supposed to establish
$\|\bar x^{(t+1)}-x^{(\star,t+1)}\|\leq\delta_t$ inserts the retracted
$x^{(t+1)}$ inside the objective gap.  APGD controls the gap of
$\bar x^{(t+1)}$, not the later retracted point.  Replacing that one symbol by
$\bar x^{(t+1)}$ yields exactly the first inequality in
\cref{eq:retraction-properties} and repairs the proof.
\end{enumerate}

Without these repairs the displayed derivation is invalid.  With them, the
later support and progress arguments use the correct quantities.

\subsection{Nonfatal algebra and notation defects}

\begin{enumerate}[leftmargin=*]
\item The PageRank off-diagonal entry is printed with denominator
$2d_id_j$ rather than $2\sqrt{d_id_j}$.  Only its nonpositive sign is used.
\item A restricted-gradient Lipschitz estimate uses $L$ where the squared
norm bound requires $L^2$.  The next logarithmic expression already contains
$L^2$, so the final soft-order statement is unchanged.
\item The volume definition alternates between a generic set and
$S^\star$, and the graph formula counts an undirected internal edge once
while $\operatorname{nnz}(Q_{S,S})$ counts both symmetric entries.  These
differ only by a factor at most two in the stated work results.
\item The complexity proof places a stage-dependent
$\widehat\varepsilon_t$ under a maximum without taking the corresponding
minimum tolerance.  Bounding it by the smallest stage tolerance repairs the
uniform logarithm.
\item The PageRank norm estimate should use coordinatewise absolute-value
bounds, replacing a displayed difference by a sum.  This changes only the
constant inside a logarithm.
\end{enumerate}

\subsection{Executable access-model obligations}

The source uses exact tests of $\nabla_iG(x)<0$ and exact active-set equality.
A numerical implementation needs sign tolerances and a proof that they do not
admit bad coordinates or terminate prematurely.  Likewise, computing the
outside gradient by accumulating active columns in a hash table takes memory
proportional to the current boundary, not $O(|S^\star|)$.  The simultaneous
soft-linear-time and $O(|S^\star|)$-space claim can be realized if adjacency
lists are sorted: merge the active neighbor streams with a heap of size
$O(|S|)$, aggregate one outside coordinate at a time, and retain only the
newly negative indices.  This adds a $\log|S|$ time factor hidden by
$\softO$, but it is an additional data-structure assumption rather than a
consequence of the pseudocode alone.

\section{Official implementation and empirical-speed audit}
\label{sec:implementation-audit}

\subsection{Provenance and what the plots actually show}

The public repository
\url{https://github.com/ZIB-IOL/AAPPR.jl} was committed in December 2023,
after the March 2023 arXiv submission and the COLT version.  Commit
\texttt{f016c46} contains the Julia implementation, serialized results, and
plots; commit \texttt{3a169eb} adds the last journal-version experiment.  The
benchmark driver uses four SNAP graphs, selects a minimum-degree vertex as the
seed, fixes $\varepsilon=10^{-6}$, sweeps either $\alpha$ or $\rho$, performs
one compilation warm-up, and records one timed run per point.  There are no
repeated seeds, error bars, or machine-independent work counts.

The official performance plots do \emph{not} support the statement that
default \ASPR{} is empirically slow relative to the repository's FISTA and
ISTA baselines.  Across the displayed sweeps, default \ASPR{} is usually
faster than both; CASPR is usually fastest.  The accurate empirical statement
is narrower:
\begin{enumerate}[leftmargin=*]
\item \ASPR{} is slower than CASPR in those experiments;
\item periodic-gradient variants
      \texttt{ASPR 1}, \texttt{ASPR 5}, and \texttt{ASPR 15}
      are slower than default \ASPR{} on the displayed Amazon and
      \texttt{ca-Astroph} sweeps; and
\item an independently implemented ASPR can still be slow for the
      theorem-level repeated-solve reason proved later in this note.
\end{enumerate}

The second observation is not evidence against early discovery as a
mathematical idea.  It is explained by a concrete implementation defect.

\subsection{The periodic boundary check is inert on RPPR}

Let $S$ be the current good set and $B=N(S)\setminus S$.  In
\path{accelerated_projected_gradient_descent}, the implementation copies
the neighborhood objective and then calls
\path{set_nonindex_entries_to_zero!} with indices $B$.  That helper
retains only the principal block $Q_{B,B}$, deleting the cross block
$Q_{B,S}$.  The alleged outside gradient is then evaluated at an iterate
$y$ supported on $S$.

\begin{proposition}[The official early-discovery flag cannot fire on RPPR]
\label{prop:official-early-inert}
Assume all initially negative constant-gradient coordinates were placed in
the first good set, as in the official implementation.  At every later
periodic check, the code computes a nonnegative value for every coordinate in
$B$.  Consequently its early-discovery flag is always false.
\end{proposition}

\begin{proof}
After the cross block is deleted and because $y_B=0$, the implemented value at
$i\in B$ is
\[
 (Q_{B,B}y_B)_i+
 \alpha\rho\sqrt{d_i}-\frac{\alpha s_i}{\sqrt{d_i}}
 =
 \alpha\rho\sqrt{d_i}-\frac{\alpha s_i}{\sqrt{d_i}}.
\]
Every coordinate where this constant is negative was selected at
initialization.  The good sets only grow, so a remaining outside coordinate
has nonnegative constant.  Hence the periodic check finds no negative outside
entry and cannot trigger its early-return branch.
\end{proof}

Thus the periodic variants perform the default inner iterations plus
additional matrix copies, restrictions, and boundary-gradient calls.  Their
slowdown in the official plots is the expected consequence of
\cref{prop:official-early-inert}, not evidence that correctly implemented
early discovery is intrinsically slower.

The correct local boundary computation is
\[
 \nabla_B G(y)=Q_{B,S}y_S-c_B.
\]
It needs the cross block and can be maintained incrementally from newly
changed active coordinates; constructing a principal outside objective is
both mathematically wrong and unnecessarily expensive.

\subsection{The implementation's early predicate is nevertheless safe}

The code checks negativity only on currently positive active entries, whereas
the source lemma assumes nonpositive gradient on all of $S$.  For a
Stieltjes quadratic, the code's weaker-looking predicate is in fact
sufficient.

\begin{lemma}[Support-restricted negativity implies a lower point]
\label{lem:positive-support-predicate}
Let $z^S>0$ on $S$ be the exact restricted minimizer.  Let $y\geq0$ be
supported on $S$, put $P:=\supp(y)$, and assume
$\nabla_PG(y)\leq0$.  Then $y\leq z^S$.  Therefore any outside coordinate
with negative gradient at $y$ belongs to the global optimal support.
\end{lemma}

\begin{proof}
Let $Z:=S\setminus P$.  Restricted optimality gives
$\nabla_SG(z^S)=0$.  On $P$,
\[
 H_{P,P}(y_P-z_P^S)
 =\nabla_PG(y)+H_{P,Z}z_Z^S\leq0,
\]
because $H_{P,Z}\leq0$.  Principal Stieltjes inverses are entrywise
nonnegative, so $y_P\leq z_P^S$; on $Z$, $y_Z=0<z_Z^S$.
The outside-gradient comparison and the support-safety argument from
\cref{lem:safe-progress} finish the proof.
\end{proof}

Accordingly, the flaw in the periodic option is the deleted cross block, not
the positive-support predicate.

\subsection{Further executable defects and benchmark caveats}

\begin{enumerate}[leftmargin=*]
\item The in-place retraction first copies the old vector, subtracts
      $\delta$ from the positive entries of the original, and then tests the
      \emph{old copy} for negativity.  If
      $0<x_i<\delta$, the result is $x_i-\delta<0$ rather than zero.  This
      violates feasibility and the lower-envelope proof.  The tests cover
      $x_i=\delta$ but not the strict interval $(0,\delta)$.
\item The prescribed APGD count is converted directly to an integer without
      the positive-part clamp in \cref{eq:corrected-inner-count}.
\item The official inner APGD also exits when its restricted gradient norm
      falls below $\mu\delta$.  Under the positive restricted-minimizer
      invariant this is a valid sufficient distance test, but it is not the
      fixed-count Algorithm~4 analyzed by
      \cref{thm:aspr-path-work}.  The literal lower bound therefore applies to
      the paper algorithm, not automatically to every run of this code.
\item The in-place support filter constructs
      \texttt{setdiff(1:n, indices)}.  ISTA and FISTA call it in every outer
      iteration, introducing an $O(n)$ allocation that is absent from the
      stated local work model and biases their wall-time comparison.
\item Neighborhood construction repeatedly concatenates, sorts, and
      deduplicates index arrays, while restricted objectives are repeatedly
      copied and mutated.  These costs can dominate sparse arithmetic for
      small active sets.
\item The algorithms do not stop by evaluating the claimed objective gap.
      They use algorithm-dependent gradient-norm logic calibrated by
      $\varepsilon$, and the diagnostic plots subtract the minimum objective
      observed among the compared runs rather than a certified optimum.
      Hence the wall times are not a fully verified equal-achieved-accuracy
      comparison.
\end{enumerate}

\section{Endpoint paths force one coordinate per stage}
\label{sec:path-expansion}

Let $P_s=(v_0,v_1,\ldots,v_{s-1})$ be an unweighted path, take the endpoint
seed $s=e_{v_0}$, and use the shared RPPR objective.  Define
\begin{equation}
 \chi:=\frac{1-\sqrt\alpha}{1+\sqrt\alpha}.
 \label{eq:chi}
\end{equation}
Let $x^\star(0)$ denote the unregularized solution and put
$y_i^\star(0):=x_i^\star(0)/\sqrt{d_i}$.  Direct solution of the path
recurrence gives
\begin{equation}
 y_i^\star(0)
 =\sqrt\alpha\,
   \frac{\chi^i+\chi^{2(s-1)-i}}{1-\chi^{2(s-1)}}
 \geq\sqrt\alpha\,\chi^i.
 \label{eq:path-green}
\end{equation}
Moreover, $Q(D^{1/2}\one)=\alpha D^{1/2}\one$.  The Stieltjes comparison
principle applied to
$[x^\star(0)-\rho D^{1/2}\one]_+$ therefore gives
\begin{equation}
 \frac{x_i^\star(\rho)}{\sqrt{d_i}}
 \geq y_i^\star(0)-\rho
 \geq \sqrt\alpha\,\chi^i-\rho.
 \label{eq:path-positive-lower}
\end{equation}
Consequently, the sufficient condition
\begin{equation}
 0<\rho<\frac12\sqrt\alpha\,\chi^{s-1}
 \label{eq:full-support-condition}
\end{equation}
ensures $S^\star(\rho)=V(P_s)$.

\begin{lemma}[Unique negative frontier]
\label{lem:unique-frontier}
Let $U_r:=\{v_0,\ldots,v_{r-1}\}$ for $1\leq r<s$, and let $z^{(r)}$ be the
exact minimizer of $G_\rho$ on the nonnegative face supported on $U_r$.
Then
\begin{equation}
 \nabla_{v_r}G_\rho(z^{(r)})<0,
 \qquad
 \nabla_{v_j}G_\rho(z^{(r)})>0\quad(j\geq r+1).
 \label{eq:unique-frontier}
\end{equation}
\end{lemma}

\begin{proof}
Every vertex $v_j$ with $j\geq r+1$ has no neighbor in $U_r$ and is not the
seed.  Hence its quadratic contribution is zero and its obstacle gradient is
$\alpha\rho\sqrt{d_j}>0$.  The only remaining outside coordinate is $v_r$.
If its gradient were nonnegative, $z^{(r)}$ would satisfy all global KKT
conditions: the restricted KKT conditions hold on $U_r$, the frontier would
be nonnegative, and all farther gradients are strictly positive.  Strong
convexity would give $z^{(r)}=x^\star(\rho)$, contradicting
$x_{v_r}^\star(\rho)>0$ from \cref{eq:full-support-condition}.
\end{proof}

Starting with $U_1$, the restricted minimizer is strictly positive.  The
negative frontier in \cref{lem:unique-frontier}, together with inverse
positivity of the enlarged principal Stieltjes matrix, makes the next
restricted minimizer strictly positive on all of $U_{r+1}$ and no smaller on
$U_r$.  Induction therefore verifies the positivity invariant required in
\cref{lem:retraction} and justifies the strict positivity of $a_*$ below.

The strict inequalities allow a quantitative finite-accuracy statement.
Define
\begin{equation}
 a_*:=\min_{1\leq r\leq s}z_{v_{r-1}}^{(r)}>0,
 \qquad
 m_*:=\min_{1\leq r<s}
       \bigl(-\nabla_{v_r}G_\rho(z^{(r)})\bigr)>0,
 \label{eq:path-margins}
\end{equation}
and let
\begin{equation}
 q_*:=\max_{\{i,j\}\in E(P_s)}|Q_{ij}|<1.
 \label{eq:path-coupling}
\end{equation}

\begin{theorem}[Exactly $s$ ASPR calls on a path]
\label{thm:s-calls}
Assume \cref{eq:full-support-condition}.  There exists
$\varepsilon_0=\varepsilon_0(\alpha,\rho,s)>0$ such that, for every
$0<\varepsilon_{\rm obj}\leq\varepsilon_0$, literal \ASPR{} generates
\begin{equation}
 S^{(r-1)}=U_r,
 \qquad r=1,\ldots,s,
 \label{eq:prefix-sequence}
\end{equation}
and therefore calls APGD exactly $s$ times before its active set stabilizes.
\end{theorem}

\begin{proof}
At zero, only the seed has negative obstacle gradient, so the first active set
is $U_1$.  At a stage with active prefix $U_r$, Lemma~\ref{lem:retraction}
gives
\begin{equation}
 \|x^+-z^{(r)}\|_2
 \leq\|x^+-\bar x\|_2+\|\bar x-z^{(r)}\|_2
 \leq(1+\sqrt r)\delta_{U_r}.
 \label{eq:path-retraction-distance}
\end{equation}
More sharply, its last coordinate differs from $z_{v_{r-1}}^{(r)}$ by at
most $2\delta_{U_r}$.  Choose $\varepsilon_0$ small enough that, for every
$r\leq s$,
\begin{equation}
 2\delta_{U_r}\leq\frac{a_*}{2},
 \qquad
 2q_*\delta_{U_r}\leq\frac{m_*}{2}.
 \label{eq:epsilon-small-enough}
\end{equation}
Then the last prefix coordinate remains positive.  The frontier-gradient
perturbation uses only the last path edge, so
\[
 \nabla_{v_r}G_\rho(x^+)
 \leq -m_*+2q_*\delta_{U_r}
 \leq-\frac{m_*}{2}<0.
\]
All vertices beyond $v_r$ retain the strictly positive regularization
gradient because they have no neighbor in $U_r$.  Thus the only new index is
$v_r$, and induction gives \cref{eq:prefix-sequence}.  At $U_s=V(P_s)$ there
is no outside coordinate, so the set stabilizes after the $s$th call.
\end{proof}

The theorem is stronger than the generic counting argument ``at most
$|S^\star|$ stages'': it gives an RPPR family on which equality holds.

\begin{proposition}[Early discovery cannot skip path layers]
\label{prop:early-path-layers}
Consider any active-set method on the same path that currently uses the prefix
$U_r$, queries gradients only at points supported on $U_r$, and admits only
negative-gradient coordinates.  A gradient query can admit at most the single
vertex $v_r$.  Thus even the source's optional early-discovery modification
requires at least $s-1$ successful expansion events.  If each event uses a
fresh incident-gradient scan of the current prefix, its discovery work is
$\Omega(s^2)$.
\end{proposition}

\begin{proof}
For every $j\geq r+1$, vertex $v_j$ is neither the seed nor adjacent to
$U_r$.  At every point supported on $U_r$,
\[
 \nabla_{v_j}G_\rho(x)=\alpha\rho\sqrt{d_j}>0.
\]
Only $v_r$ can have negative outside gradient.  Hence at most one path layer
is admitted per successful query.  A fresh incident scan at prefix length
$r$ costs $\Theta(r)$, and summing over $r=1,\ldots,s-1$ gives the result.
\end{proof}

This proposition separates two effects.  Correct early discovery may avoid
solving every prefix to the final inner tolerance, so the
$1/\sqrt\alpha$ factor in the literal-Algorithm~4 lower bound need not survive
unchanged.  It cannot, however, remove the $s$ sequential path discoveries;
with fresh scans it still pays quadratic discovery work.  An implementation
that incrementally maintains only the frontier residual is outside this scan
model and can avoid that particular $\Omega(s^2)$ charge.

\section{A matching work lower bound for literal ASPR}
\label{sec:work-lower-bound}

At stage $r+1$, the starting point is supported on $U_r$ and has zero entry
at the newly admitted coordinate $v_r$.  The preceding proof gives
\begin{equation}
 \|\nabla_{U_{r+1}}G_\rho(x^{(r)})\|_2
 \geq\frac{m_*}{2}
 \qquad(1\leq r<s).
 \label{eq:starting-gradient-lower}
\end{equation}
The first call has the fixed nonzero seed gradient.  Let $g_*>0$ be the
minimum of these finitely many lower bounds.  Substituting
\cref{eq:inner-tolerances} into the inner-count ratio gives, for RPPR,
\begin{equation}
 \frac{(1-\alpha)\|\nabla_{U_r}G_\rho(x^{(r-1)})\|_2^2}
      {2\widehat\varepsilon_{U_r}\alpha^2}
 \geq
 \frac{(1-\alpha)(1+r)g_*^2}
      {\varepsilon_{\rm obj}\alpha^4}.
 \label{eq:inner-ratio-lower}
\end{equation}
After decreasing $\varepsilon_0$ once more, the right-hand side is at least
$e$.  The corrected prescribed count \cref{eq:corrected-inner-count} then
satisfies
\begin{equation}
 K_{U_r}\geq 2\sqrt{1/\alpha}
 \qquad(r=1,\ldots,s).
 \label{eq:inner-count-lower}
\end{equation}

\begin{theorem}[ASPR path-work lower bound]
\label{thm:aspr-path-work}
Under the assumptions of Theorem~\ref{thm:s-calls}, for all sufficiently
small $\varepsilon_{\rm obj}$ the literal \ASPR{} algorithm satisfies
\begin{equation}
 \Work_{\rm inner}
 \geq c\frac{s^2}{\sqrt\alpha}
 \label{eq:inner-work-lower}
\end{equation}
for a universal work-model constant $c>0$.  If every expansion line computes
the full gradient afresh from the active columns, then
\begin{equation}
 \Work_{\rm discover}\geq c' s^2.
 \label{eq:discovery-work-lower}
\end{equation}
\end{theorem}

\begin{proof}
A path prefix of length $r$ has $r$ diagonal entries and $2(r-1)$ symmetric
off-diagonal entries in its principal matrix.  Thus a restricted gradient
touches $3r-2=\Theta(r)$ stored entries; even under the source's once-per-edge
internal-volume convention, the cost is $2r-1=\Theta(r)$.  Combining this
with \cref{eq:inner-count-lower},
\[
 \Work_{\rm inner}
 \geq c_0\sum_{r=1}^s\frac{r}{\sqrt\alpha}
 =\Omega\!\left(\frac{s^2}{\sqrt\alpha}\right).
\]
A fresh incident-column scan at stage $r$ touches $\Theta(r)$ path entries,
so summing the $s$ scans gives \cref{eq:discovery-work-lower}.
\end{proof}

\begin{corollary}[A same-tolerance factor-$s$ separation]
\label{cor:quantitative-separation}
For integers $s\geq2$, set
\begin{equation}
 \alpha=s^{-2},
 \qquad
 \rho=\frac{\alpha}{100},
 \qquad
 \varepsilon_{\rm obj}=10^{-4}\alpha^2.
 \label{eq:quantitative-family}
\end{equation}
Then this endpoint-path family has full optimal support and
\begin{equation}
 \Work_{\ASPR}=\Omega(s^3\log s),
 \qquad
 \Work_{\FISTA}=O(s^2\log s).
 \label{eq:quantitative-separation}
\end{equation}
Thus the same instances at the same objective tolerance exhibit a factor
$\Omega(s)$ separation.
\end{corollary}

\begin{proof}
The full-support condition follows from
$\chi=(s-1)/(s+1)$ and
$\chi^{s-1}\geq e^{-2}$.  To quantify the finite-accuracy margins, write
\[
 b:=\frac{1-\alpha}{2},
 \qquad
 \theta:=\log\frac{1+\sqrt\alpha}{1-\sqrt\alpha}.
\]
For an unregularized proper prefix $U_r$, the degree-normalized last
coordinate is obtained by solving the Dirichlet path recurrence:
\begin{equation}
 \frac{z_{v_{r-1}}^{(r,0)}}{\sqrt{d_{r-1}}}
 =\frac{\alpha}{b\cosh(r\theta)}.
 \label{eq:prefix-green-last}
\end{equation}
Here $r\theta\leq s\theta\leq2\log3$, so the right-hand side is at least
$2\alpha/5$.  If $Q_r:=Q_{U_r,U_r}$ and
$w_r:=(D^{1/2}\one)_{U_r}$, then
\[
 Q_rw_r\geq\alpha w_r,
 \qquad
 Q_r^{-1}\alpha w_r\leq w_r.
\]
Consequently the regularized restricted solution is bounded below by the
unregularized one minus $\rho w_r$.  Every proper-prefix last coordinate is
therefore at least $39\alpha/100$ after degree normalization.  The full-path
last coordinate has the same order by \cref{eq:path-green}.  In particular,
\[
 a_*\geq\frac{\alpha}{5}.
\]
For a proper prefix, multiplying the frontier gradient by
$\sqrt{d_r}$ gives
\[
 \sqrt{d_r}\,
 \bigl(-\nabla_{v_r}G_\rho(z^{(r)})\bigr)
 =b\,\frac{z_{v_{r-1}}^{(r)}}{\sqrt{d_{r-1}}}
  -\alpha\rho d_r.
\]
Since $b\geq3/8$ and $d_r\leq2$, this yields
$m_*\geq\alpha/20$.

Under \cref{eq:quantitative-family},
$\delta_{U_r}\leq10^{-2}\alpha^{3/2}/\sqrt2$.  The two strict-margin
requirements in \cref{eq:epsilon-small-enough} hold, so
\cref{thm:s-calls} applies.  The seed gradient is at least $\alpha/2$, and
every later starting gradient is at least $m_*/2$.  Thus the ratio in
\cref{eq:inner-ratio-lower} is $\Omega(\alpha^{-4})$, and the prescribed count
at every stage is
\[
 K_{U_r}=\Omega\!\left(
   \alpha^{-1/2}\log(1/\alpha)
 \right)=\Omega(s\log s).
\]
Summing the $\Theta(r)$ restricted-gradient costs proves the \ASPR{} bound.

For FISTA, $\|c_\rho\|_2^2=O(\alpha^2)$ on this family, and hence
$\Delta_0=G_\rho(0)-G_\rho(x^\star)\leq
\|c_\rho\|_2^2/(2\alpha)=O(\alpha)$.  Its iteration count is therefore
$O(\alpha^{-1/2}\log(\Delta_0/\varepsilon_{\rm obj}))
=O(s\log s)$.  Each full-path gradient costs $\Theta(s)$, proving the second
bound in \cref{eq:quantitative-separation}.
\end{proof}

On this family,
\begin{equation}
 |S^\star|=s,
 \qquad
 \widetilde\vol(S^\star)=\Theta(s),
 \qquad
 \vol(S^\star)=\Theta(s),
 \qquad L=1.
 \label{eq:path-parameters}
\end{equation}
Therefore Theorem~\ref{thm:aspr-path-work} matches both structural products in
the published upper bound:
\begin{equation}
 |S^\star|\widetilde\vol(S^\star)\sqrt{L/\alpha}
 =\Theta(s^2/\sqrt\alpha),
 \qquad
 |S^\star|\vol(S^\star)=\Theta(s^2),
 \label{eq:matching-products}
\end{equation}
up to the logarithm in the prescribed APGD accuracy.  This establishes
algorithm-specific tightness even within normalized RPPR, not merely within
the larger class of arbitrary M-matrix quadratics.

\section{Why standard FISTA can be linearly faster}
\label{sec:fista-comparison}

Run standard strongly convex \FISTA{} directly on the same finite path
problem.  Regardless of transient sparsity, every iterate is supported in a
graph of total degree volume $\Theta(s)$.  The standard objective-gap rate
uses
\begin{equation}
 N_{\varepsilon}
 =O\!\left(
  \frac1{\sqrt\alpha}
  \log\frac{\Delta_0}{\varepsilon_{\rm obj}}
 \right)
 \label{eq:fista-iterations}
\end{equation}
iterations.  Hence
\begin{equation}
 \Work_{\FISTA}
 =O\!\left(
  \frac{s}{\sqrt\alpha}
  \log\frac{\Delta_0}{\varepsilon_{\rm obj}}
 \right).
 \label{eq:fista-path-work}
\end{equation}
The structural comparison is a factor $s$ before accuracy logarithms.
\Cref{cor:quantitative-separation} shows that this is also a literal
same-tolerance factor-$s$ separation, rather than an artifact of suppressing a
path-length-dependent logarithm.

This does not contradict the star lower bound of
\citet{fountoulakis2026}.  On their leaf-seeded star, standard FISTA
transiently activates a center of arbitrarily large degree even though the
optimum is seed-supported, whereas \ASPR{} stays in the true support.  The two
families establish incomparability:
\begin{itemize}[leftmargin=*]
\item a full-support path punishes repeated active-set restarts and favors
      continuation by standard FISTA;
\item a tiny-support, high-degree-boundary star punishes momentum overshoot
      and favors ASPR's safe support containment.
\end{itemize}

\section{What is and is not tight}
\label{sec:tightness-scope}

The path theorem closes the main algorithmic question:
\begin{enumerate}[leftmargin=*]
\item The maximum number of APGD calls is exactly $|S^\star|$ in the worst
      case.
\item Repeated restricted-gradient work can attain
      $|S^\star|\widetilde\vol(S^\star)/\sqrt\alpha$ up to logarithms.
\item Literal recomputation of every discovery gradient can attain
      $|S^\star|\vol(S^\star)$.
\end{enumerate}

Four stronger interpretations remain false or open.
\begin{enumerate}[leftmargin=*]
\item The result is not a lower bound for an implementation that preserves an
      accelerated estimate sequence across expansions.  Such a method is not
      Algorithm~4.
\item The optional early-discovery variant still needs $|S^\star|$ successful
      events and $\Omega(s^2)$ fresh-scan work on the path, but the present
      proof does not force it to spend $\Omega(1/\sqrt\alpha)$ inner
      iterations before every event.
\item It is not an information-theoretic lower bound for local RPPR.  Standard
      FISTA already beats ASPR on the path, while safe local continuation may
      do better still.
\item The sufficiently-small threshold
      $\varepsilon_0(\alpha,\rho,s)$ is an objective-gap threshold.  A fair
      lower bound under APPR accuracy, PPR infinity error, or a proximal
      fixed-point residual requires an explicit conversion that the current
      repository has not adopted.
\end{enumerate}

The lower bound identifies the precise design target for a successor to
ASPR.  It must retain the safe lower-envelope or retraction idea without
restarting a high-accuracy accelerated solve after every frontier event.  A
useful theorem would charge
\begin{equation}
 \begin{aligned}
  \sum_t K_t\widetilde\vol(S_t)
  &\quad\text{by accelerated continuation, and}\\
  \sum_t\bigl(\vol(S_{t+1})-\vol(S_t)\bigr)
  &\quad\text{by newly admitted volume},
 \end{aligned}
 \label{eq:desired-amortization}
\end{equation}
rather than multiplying both final volumes by $|S^\star|$.

The new safe-predicate lemma isolates the remaining continuation problem.
Under the positive-restricted-minimizer invariant, at an accelerated iterate
$y$ one may safely expand whenever
\[
 \nabla_{\supp(y)}G(y)\leq0
 \quad\text{and}\quad
 \min_{j\in N(S)\setminus S}\nabla_jG(y)<0,
\]
provided the boundary gradient retains the cross term $Q_{B,S}y_S$.  No
high-accuracy restricted solution is needed for that support-safety step.
What is still missing is an expanding-face estimate-sequence lemma showing
that the accelerated potential survives embedding its momentum state into
$S^+\supset S$.  Without such a lemma, carrying momentum is a plausible
algorithmic proposal but not a proved accelerated rate.  This, rather than
support safety, is the precise theorem that a continuation-based hybrid must
establish.

\section{Conclusion}

The COLT 2023 ASPR theorem has a correct mathematical core but needs an
explicit repaired formulation.  More importantly, its pessimistic-looking
$|S^\star|$ multiplier is a real property of the algorithm.  Endpoint paths
force one coordinate per stage, and sufficiently accurate nested solves force
the full $\Omega(s^2/\sqrt\alpha)$ work.  The official later experiments do
not show default \ASPR{} losing to their FISTA baseline, so they must not be
cited as evidence of empirical ASPR slowness.  They do show periodic-scan
variants losing to default \ASPR{}, which
\cref{prop:official-early-inert} explains as an implementation defect.
Independent slow implementations can still realize the repeated-prefix
worst case proved here.  The safe-support mechanism remains valuable, but it
should be combined with continuation and incremental frontier maintenance
rather than nested restarts and repeated scans.

\footnotesize
\begin{thebibliography}{9}

\bibitem[Mart\'inez-Rubio et~al.(2023)Mart\'inez-Rubio, Wirth, and
Pokutta]{martinezrubio2023}
David Mart\'inez-Rubio, Elias Wirth, and Sebastian Pokutta.
\newblock Accelerated and sparse algorithms for approximate personalized
PageRank and beyond.
\newblock In \emph{Proceedings of the 36th Conference on Learning Theory},
pages 2852--2876, 2023.

\bibitem[Fountoulakis and Mart\'inez-Rubio(2026)]{fountoulakis2026}
Kimon Fountoulakis and David Mart\'inez-Rubio.
\newblock Complexity of classical acceleration for
$\ell_1$-regularized PageRank.
\newblock arXiv:2602.21138v2, 2026.

\bibitem[Diakonikolas and Orecchia(2019)]{diakonikolas2019}
Jelena Diakonikolas and Lorenzo Orecchia.
\newblock The approximate duality gap technique: A unified theory of
first-order methods.
\newblock \emph{SIAM Journal on Optimization}, 29(1):660--689, 2019.

\bibitem[Fountoulakis et~al.(2019)Fountoulakis, Roosta-Khorasani,
Shun, Cheng, and Mahoney]{fountoulakis2019}
Kimon Fountoulakis, Farbod Roosta-Khorasani, Julian Shun, Xiang Cheng, and
Michael W. Mahoney.
\newblock Variational perspective on local graph clustering.
\newblock \emph{Mathematical Programming}, 174:553--573, 2019.

\end{thebibliography}

\end{document}
